%% Part 2.2 resolution fragment: v1.1 Part III, lines 5568--5958.
%% The two inference rules live only on registry sheet resolution-fol.

\chapter{First-Order Clauses, Unification, and Resolution}
\label{sec:p02-first-order-resolution}
\chaptermark{First-Order Resolution}

\section{Clauses and simultaneous term substitution}
\label{sec:p02-fol-clauses}

\begin{hbdefinition}{def:p02-fol-clause}{First-order clause and formula interpretation}
A \emph{first-order clause}\index{clause!first-order} is a finite set of
first-order literals; order and multiplicity are not clause data.  To obtain
one exact formula presentation, supply a well-order of the signature's symbol
tokens, combine it with the variable and logical-token orders, and order
formula strings first by length and then lexicographically.  If the resulting
listing of a nonempty clause \(C\) is \(L_1\prec\cdots\prec L_k\), define
\(\operatorname{form}_{\mathcal L}(C)\) as their right-associated
disjunction; define the empty clause's formula to be \(\bottom\).

If \(x_{i_1},\ldots,x_{i_m}\), in increasing index order, are exactly the
variables occurring in \(C\), put
\[
 \operatorname{ucl}_{\mathcal L}(C)
 \defeq
 \forall x_{i_1}\cdots\forall x_{i_m}
 \operatorname{form}_{\mathcal L}(C).
\]
Both maps are metalinguistic.  A structure satisfies \(C\) iff it satisfies
this \emph{universal closure}\index{universal closure!of a clause}, and
satisfies a clause set iff it satisfies every
member.  The supplied well-order affects presentation, not satisfaction; for
an arbitrary signature its existence is the explicit well-ordering/choice
assumption.

A clause is \emph{equality-free}\index{clause!equality-free} iff none of its
literals contains logical
identity.  This restricts clause contents, not the signature, since equality
is logical rather than a signature member.  A clause set is equality-free iff
each of its members is equality-free.
\end{hbdefinition}

\begin{hbdefinition}{def:p02-fol-simultaneous-substitution}{Finite-support simultaneous term substitution}
A \emph{term substitution}\index{substitution!simultaneous term} is a total map
\(\theta:\mathsf{Var}\to\operatorname{Term}(\mathcal L)\) whose
\emph{variable support}\index{substitution support!variable support}
\[
 \operatorname{supp}_{\mathrm{var}}(\theta)
 \defeq\{x:\theta(x)\ne x\}
\]
is finite.  This support is a set of variables, unlike the multiplicity
support of a finite multiset.  A finite display
\(\{x_1\mapsto t_1,\ldots,x_n\mapsto t_n\}\) maps the listed variables as
shown and fixes every other variable.  Extend \(\theta\) homomorphically to
terms, atoms, literals, and clauses, writing
\(t\theta,A\theta,L\theta,C\theta\).
The action is simultaneous: it does not reapply \(\theta\) inside an inserted
term.  On clauses it takes the pointwise image, so set equality removes
duplicates.  It is a quantifier-free operation distinct from the earlier
total capture-avoiding single-variable formula substitution.
\end{hbdefinition}

\begin{hbdefinition}{def:p02-fol-instance-preorder}{Composition and instance preorder}
Let \(\mathrm{id}\) fix every variable, and orient composition by
\[
 (\theta\eta)(x)\defeq(\theta(x))\eta.
\]
Then
\(\operatorname{supp}_{\mathrm{var}}(\theta\eta)\subseteq
  \operatorname{supp}_{\mathrm{var}}(\theta)
  \cup\operatorname{supp}_{\mathrm{var}}(\eta)\),
so the composite again has finite support.
Define the \emph{instance preorder}\index{instance preorder} by
\[
 \theta\preceq_{\mathrm{inst}}\rho
 \quad\text{iff}\quad
 (\exists\eta)\;\rho=\theta\eta.
\]
This metalinguistic relation compares substitutions by further
instantiation; it is not equality.  Identity and associative composition make
it reflexive and transitive.  Its later use is exactly the definition of
``most general.''
\end{hbdefinition}

\begin{hblemma}{thm:p02-fol-substitution-composition}{proved=annex:p2.fol.res.1}{Composition of simultaneous substitutions}
For every term \(t\), atom \(A\), literal \(L\), and clause \(C\),
\[
 t(\theta\eta)=(t\theta)\eta,
 \quad A(\theta\eta)=(A\theta)\eta,
 \quad L(\theta\eta)=(L\theta)\eta,
 \quad C(\theta\eta)=(C\theta)\eta.
\]
Consequently substitution composition is associative.
\end{hblemma}

\begin{hblemma}{thm:p02-fol-clause-substitution}{proved=annex:p2.fol.res.2}{Satisfaction under simultaneous term substitution}
If an \(\mathcal L\)-structure \(\mathcal M\) satisfies a clause \(C\),
then it satisfies \(C\theta\) for every finite-support term substitution
\(\theta\).
\end{hblemma}

\section{Unification}
\label{sec:p02-fol-unification}

\begin{hbdefinition}{def:p02-fol-renaming-apart}{Renaming apart}
A \emph{variable renaming}\index{variable renaming} is a term substitution
with finite support, induced by
a bijection \(\pi:\mathsf{Var}\to\mathsf{Var}\).  Put
\(\operatorname{Var}(C)\defeq
  \bigcup_{L\in C}\operatorname{Var}(L)\).
Two clause occurrences \(C\pi,D\rho\) are
\emph{renamed apart}\index{renaming apart} iff
\[
 \operatorname{Var}(C\pi)\cap\operatorname{Var}(D\rho)=\varnothing.
\]
Such renamings exist because each
clause uses finitely many variables and the variable supply is infinite.
Universal closure makes the renamed clause equivalent to its original.
This operation is required before every binary inference on
\hyperref[sheet:resolution-fol]{\texttt{sheet:resolution-fol}}; it is not the
alpha-equivalence relation on quantified formula strings.
\end{hbdefinition}

\begin{hbdefinition}{def:p02-fol-mgu}{Unifier and most general unifier}
A term substitution \(\theta\) \emph{unifies}\index{unifier} atoms \(A,B\) iff
\(A\theta\) and \(B\theta\) are literally identical.  It is a
\emph{most general unifier}\index{most general unifier} (MGU) iff
\(\theta\preceq_{\mathrm{inst}}\rho\) for every unifier \(\rho\),
equivalently every unifier is \(\theta\eta\) for some substitution
\(\eta\).  Most general therefore refers to the instance preorder, not to
syntactic size or literal uniqueness.
\end{hbdefinition}

\begin{hbmethod}{def:p02-fol-unification-procedure}{First-order unification procedure}
\index{unification!first-order procedure}
First fail if the two atoms have different outer predicate symbols or arities;
for equality atoms, regard equality as the common binary outer symbol and pair
their two ordered arguments.  Otherwise form the finite set of equations
between corresponding argument terms and initialize the accumulated
substitution to \(\mathrm{id}\).  Repeatedly delete \(t\doteq t\), decompose
matching function symbols, fail on a symbol/arity clash, orient a variable to
the left, or eliminate \(x\doteq t\): if \(x\) occurs properly in \(t\),
fail; otherwise put \(\delta=\{x\mapsto t\}\), apply \(\delta\)
simultaneously to every remaining equation, and replace the accumulator
\(\theta\) by \(\theta\delta\).  Return the accumulator when
no equation remains.  The occurs check is a procedure step needed because
finite terms cannot contain themselves; it is not part of the definition of
a substitution.
\end{hbmethod}

\begin{hbtheorem}{thm:p02-fol-unification}{sketched=annex:p2.fol.res.3}{First-order unification}
The procedure terminates.  It returns failure iff the input atoms are not
unifiable; otherwise it returns a most general unifier.
\end{hbtheorem}

\section{Resolution and factoring}
\label{sec:p02-fol-resolution-calculus}

\begin{hbdefinition}{def:p02-fol-resolution-derivation}{First-order resolution derivation}
A \emph{first-order resolution derivation from \(S\)}
\index{resolution!first-order derivation} is a finite sequence
of clauses, each an input member of \(S\) or the conclusion of a registered
Resolution or Factoring inference from earlier clauses.  Those rules are on
\hyperref[sheet:resolution-fol]{\texttt{sheet:resolution-fol}}.
Every binary inference uses renamed-apart parent
occurrences and the MGU stipulated on the sheet.  This named derivation
object is needed by the soundness and completeness statements below and is
distinct from a propositional resolution DAG.
\end{hbdefinition}

\begin{figure}[htbp]
\centering
\begin{tikzpicture}[
  font=\small,
  node distance=10mm and 18mm,
  clause/.style={draw,rounded corners,align=center,inner sep=4pt},
  >=Stealth]
  \node[clause] (c1) {\(C_1=\{P(x),Q(x)\}\)};
  \node[clause,right=of c1] (c2) {\(C_2=\{\neg P(a)\}\)};
  \node[clause,right=of c2] (c3) {\(C_3=\{\neg Q(a)\}\)};
  \node[clause,below=of c2] (c4) {\(C_4=\{Q(a)\}\)};
  \node[clause,below=of c4] (empty) {empty clause \(\varnothing\)};
  \draw[->] (c1) -- node[left,align=right,font=\footnotesize]
    {Resolution\\\(\{x\mapsto a\}\)} (c4);
  \draw[->] (c2) -- (c4);
  \draw[->] (c4) -- node[left,font=\footnotesize]{Resolution} (empty);
  \draw[->] (c3) -- (empty);
\end{tikzpicture}
\caption{Resolution-refutation DAG (apparatus A-121).  Each downward node is
the resolvent of its two incoming parents; the final node is the empty clause.
The first parents are already renamed apart because \(C_2\) is ground, and
their MGU is the displayed finite-support substitution.}
\label{fig:p02-fol-resolution-dag}
\end{figure}
\index{resolution!refutation DAG}

\begin{hbtheorem}{thm:p02-fol-resolution-soundness}{proved=annex:p2.fol.res.4}{Resolution and Factoring soundness}
Every structure satisfying the parent clauses of a registered Resolution or
Factoring inference satisfies its conclusion.  Consequently, if a
first-order resolution derivation from \(S\) contains the empty clause, then
\(S\) is unsatisfiable.
\end{hbtheorem}

\begin{hbdefinition}{def:p02-fol-ground-instance}{Ground term and ground instance}
A term \(t\) is \emph{ground}\index{ground term} iff
\(\operatorname{Var}(t)=\varnothing\).  A \emph{ground instance} of a
clause \(C\)\index{ground instance} is \(C\theta\), where \(\theta\) sends
every variable of
\(C\) to a ground term.  These notions classify terms and substitution
instances; neither means a closed formula.
\end{hbdefinition}

\begin{hbtheorem}{thm:p02-herbrand-clauses}{cited=Robinson1965}{Herbrand theorem for clauses}
Let \(S\) be a finite equality-free set of first-order clauses in a signature
with at least one ground term.  Then \(S\) is unsatisfiable iff some finite
set of ground instances of members of \(S\), regarded as propositional
clauses whose distinct ground atoms are propositional variables, is
propositionally unsatisfiable.
\end{hbtheorem}

\begin{hbtheorem}{thm:p02-resolution-lifting}{cited=BachmairGanzinger2001}{Lifting lemma}
Let renamed-apart clauses \(C,D\) have disjoint variable sets, and let \(G\)
be a ground resolvent of ground instances of them.  After factoring each
parent literal family identified by the ground substitution, there is a
first-order resolvent \(R\) obtained with an MGU and a ground substitution
\(\eta\) such that \(R\eta=G\).  The analogous assertion holds for
factoring.  Hence every finite ground resolution refutation of ground
instances of \(S\) lifts step by step to a first-order derivation with
factoring from \(S\).
\end{hbtheorem}

\begin{hbtheorem}{thm:p02-fol-resolution-completeness}{sketched=annex:p2.fol.res.5}{Refutational completeness of first-order resolution}
Let \(S\) be a finite equality-free set of first-order clauses.  After adding
a constant if the signature has no ground term, \(S\) is unsatisfiable iff
Resolution with Factoring derives the empty clause.
\end{hbtheorem}

\begin{hbwarning}{def:p02-fol-resolution-equality-boundary}{Bare resolution is not equality-complete}
The preceding completeness theorem deliberately excludes every clause
containing logical equality.  Equality atoms may be carried through the sheet
for soundness, but bare Resolution and Factoring do not provide a complete
identity treatment.  Equality axioms, paramodulation, or superposition are
additional calculi; superposition is a Phase-5 Part-12 entry.
\end{hbwarning}

The Herbrand universe/structure validity form required by D-P3-003 is stated
separately from Theorem~\ref{thm:p02-herbrand-clauses}.  The former concerns
validity of a prenex existential sentence and a finite disjunction of its
ground instances.  The latter is the dual universal-clause unsatisfiability
form used in the resolution sketch.  Neither is the quotient term model of
the Henkin construction.
Complete sources for the Herbrand clause theorem and the lifting machinery
used here are \cite{Robinson1965,BachmairGanzinger2001}.
